\subsection{Workflow Contracts}
The repository suite passed 171 tests, including 263 bounded configurations
for missingness, symptom onset, retry, correction, and report history. An
explicit-null reference agreed in 27 static F/A/S combinations. We extended the
comparison to an independent append-only event-log reference that reconstructs
measurement state, symptom activity, and urgency without invoking the production
reducer. Seeds 0--19 each generated a 30-operation pseudorandom sequence. The same 20 sequences were replayed for F/A/S, yielding 60 component runs and 1,800 transition checks. Fixed seeds make the generated interactions reproducible; the checks are not independent participant observations. Operations include measurement updates,
retries, symptom reports, and corrections. All transitions agreed in measurement
state, event activity, and urgency; all previously emitted report snapshots remained
unchanged. Measurement results are injected at the internal result boundary, so
these trajectories evaluate state management rather than sensor performance.
B/E questionnaire semantics are covered by the separate bounded configurations.


Five synthetic landmark trajectories, initialized after entry into the hold phase, exercised the hold/return implementation:
stable hold, sudden lowering, gradual lowering, brief middle occlusion, and
missing late endpoint. The implementation returned negative, positive,
positive, negative, and insufficient, respectively. Visible lowering remained
in the measurement sequence after entry into the hold phase. A regression
also verified that the passive-speech diagnostic trace associates anomaly votes
with acquisition windows. Targeted removal of symptom retention, snapshot isolation, or correction history broke the corresponding contract in each of 15 constructed cases; the full implementation satisfied all 15. Each case isolates a mechanism under a specified failure-triggering interaction. The host-only 1,000-record queue recovery test is retained in the software artifacts.
